\scanpage{67}
% archive foliation: 66

\begin{center}- 26 -\end{center}

\[
(19)\qquad
\frac{dv}{dt}=
\frac{M}{(1-\alpha t)(1-p\alpha v)(1-\alpha\sum_{\ell=1}^{q}w_\ell)}
\]
et réciproquement ; or l'intégrale de cette équation qui s'annule pour $t=0$ satisfait à l'équation
\[
v-\frac12p\alpha v^2=
-\frac{M}{\alpha(1-\alpha\sum_{\ell=1}^{q}w_\ell)}\log(1-\alpha t)
\]
d'où on tire
\[
v(t)=p\alpha\left[1-\sqrt{1+\frac{2pM}{1-\alpha\sum_{\ell=1}^{q}w_\ell}
\log(1-\alpha t)}\right]
\]
Comme $\log(1-\alpha t)$ est égale à la somme d'une série entière convergente et sans terme constant pour $|t|<\frac1\alpha$, il résulte encore du théorème de substitution (\S\,1,th.2) que $v(t)$ est égal à une série entière en $t,w_1,\ldots,w_q$ sans terme indépendant de $t$, et sommable dans un polycylindre de centre l'origine. Le th.2 est ainsi complètement démontré.

\uline{Remarque}.-- Lorsque $K=\mathbb R$, les $f_i$ sont des fonctions \uline{localement lipschitziennes} dans $A$ par rapport aux $y_j$ (\S\,1,prop.8) ; on sait alors que dans un voisinage de $a$ (pour chaque système de valeurs de $z_1,\ldots,z_q$) le système d'équations différentielles (12) n'a qu'\uline{une seule} solution satisfaisant à (13) (\uline{Fonct.var.réelle}, chap.IV,\S\,1) ; le th.2 montre que cette solution est holomorphe au voisinage de $a$.

\uline{4. Applications} : \uline{I}. \uline{Dépendance des conditions initiales}.

PROPOSITION 1.-- \uline{Soit} $K$ \uline{un corps valué complet de caractéristique} $0$.
\uline{Soient} $f_i(x;y_1,\ldots,y_p)$ $(1\leq i\leq p)$ $p$ \uline{fonctions holomorphes scalaires définies dans une partie ouverte}
$A=B\times C$ \uline{de} $K^{p+1}$ ($B$ ouvert dans $K$,$C$ ouvert dans $K^p$).
\uline{Pour tout point} $(a,b)$ \uline{de} $A$, \uline{il existe un voisinage} $U\subset B$ \uline{de} $a$ \uline{dans} $K$ \uline{et un voisinage} $V\subset C$ \uline{de} $b$ \uline{dans} $K^p$, \uline{et une fonction} $\mathbf{u}(x,x_0,\mathbf{y}_0)=(u_i(x,x_0,\mathbf{y}_0))$ \uline{à valeurs dans} $C$, \uline{et holomorphe dans} $U\times U\times V$, \uline{telle que pour tout point} $(x_0,\mathbf{y}_0)\in U\times V$, \uline{la fonction}
$x\longrightarrow \mathbf{u}(x,x_0,\mathbf{y}_0)$ \uline{soit la seule fonction holomorphe dans un voisinage de} $x_0$, \uline{vérifiant l'équation différentielle}
